\documentclass[10pt,conference,compsocconf]{IEEEtran}

\usepackage{hyperref}
\usepackage{graphicx}
\usepackage{booktabs}
\usepackage{amsmath}
\usepackage{xcolor}
\usepackage{multirow}

\addtolength{\topmargin}{-32pt}
\addtolength{\textheight}{32pt}

\begin{document}
\title{\vspace{-8pt}How Foveated Vision Shapes Cognitive Maps in Navigation Agents\vspace{-12pt}}

\author{
  Leo Bruneau, Weilun Xu (394449), Rim Abkari, Zyad Aoutir\\
  \textit{CS-503 Visual Intelligence --- Final Report}
}

\maketitle

% ==========================================================================
% ABSTRACT
% ==========================================================================
\begin{abstract}
Human spatial cognition is built on foveated vision: we see sharply only at the centre of gaze, yet construct rich cognitive maps of our surroundings. In artificial navigation agents, map-like representations emerge in recurrent memory, but always under either conditions of blind or spatially uniform perception, leaving open how non-uniform input shapes what memory encodes. We train navigation agents in Habitat under varying visual conditions---blind, uniform, foveated (fixed centre), matched-compute, and foveated (learned gaze)---and probe the resulting LSTM hidden states with linear decoders. We find that spatially non-uniform vision induces compensatory temporal memory: the foveated agent retains position information over longer horizons than all other conditions, including blind agents that receive GPS directly. Cross-condition CKA and probe-transfer analyses reveal that each visual condition produces a fundamentally different representational geometry, despite similar task performance. These results establish sensor structure---not just architecture or objective---as a causal factor in what recurrent memory learns to represent.
\end{abstract}

% ==========================================================================
% 1. INTRODUCTION
% ==========================================================================
\section{Introduction}

Wijmans et al.~\cite{wijmans2023emergence} showed that RL agents trained on PointGoal navigation develop map-like representations in their recurrent hidden states---even \emph{blind} agents receiving only goal displacement, GPS, compass, and a close-to-goal scalar. Dwivedi et al.~\cite{dwivedi2022isee} probed sighted agents for related spatial concepts. However, all these settings keep visual certainty spatially constant, leaving open whether the spatial representations are a property of the navigation task or an artifact of uniform input quality.

In biological vision, input quality is never uniform: humans see sharply only within $\sim$2\textdegree{} of gaze. Wirth et al.~\cite{wirth2017gaze} showed that primate hippocampal spatial representations depend on what the animal foveates, not just where it is---suggesting that spatial memory is shaped by how reliably an agent observes each region. If this holds for learned agents, a foveated sensor should not merely degrade cognitive maps but reshape them: memory may develop qualitatively different strategies to compensate for uneven observation quality.

Two accounts---passive buffering vs.\ active compensation---make divergent predictions. A passive memory's accuracy should fall with degraded input, since peripherally-seen regions are stored less faithfully. An uncertainty-aware memory predicts the opposite: accuracy staying flat or rising, because the agent devotes representational capacity to uncertain regions. We ask: \textbf{does a foveated agent's memory follow the passive-drop pattern, or does it exhibit the compensatory signature of epistemic representation?}

We address this question through three hypotheses. \textbf{H1 (Compensatory memory):} foveated agents build richer temporal traces to compensate for perceptual uncertainty. \textbf{H2 (Representational divergence):} foveation produces qualitatively different spatial codes, not merely degraded versions of the uniform code. \textbf{H3 (Epistemic gaze):} learned gaze targets memory-weak regions, driven by task pressure alone.

Our results support H1 and H2. The foveated agent retains spatial information over longer temporal horizons than any other condition (including blind), and cross-condition representational similarity is near zero despite comparable task performance. These findings establish that sensor structure---alongside architecture and objective---is a causal axis for the structure of learned spatial representations.

% ==========================================================================
% 2. RELATED WORK
% ==========================================================================
\section{Related Work}

\textbf{Cognitive maps in RL agents.}\quad
Beyond Wijmans et al.~\cite{wijmans2023emergence} and Dwivedi et al.~\cite{dwivedi2022isee}, Banino et al.~\cite{banino2018grid} and Cueva \& Wei~\cite{cueva2018grid} found grid-cell-like representations in recurrent networks trained on path integration, and Gornet \& Thomson~\cite{gornet2024cognitive} extended this via visual predictive coding. Hennig et al.~\cite{hennig2023belief} showed that RL under partial observability can induce belief-state-like representations without explicit probabilistic objectives---a precedent for our H1. Ramakrishnan et al.~\cite{ramakrishnan2025space} examined spatial cognition in frontier models. All these studies use uniform or absent perception; none vary input quality spatially.

\textbf{Foveated vision in deep learning.}\quad
Mnih et al.~\cite{mnih2014attention} introduced recurrent attention models. Pourrahimi \& Bashivan~\cite{pourrahimi2025foveated} probed a foveated visual-search model for neuroscience-predicted features; Shang \& Ryoo~\cite{shang2023sugarl} jointly learned gaze and motor policies. These address visual search or task performance, not how foveation reshapes a navigation agent's learned spatial memory.

\textbf{Neuroscience and sensor structure.}\quad
Wirth et al.~\cite{wirth2017gaze} established that primate hippocampal spatial representations depend on gaze, not just location. Atanov et al.~\cite{atanov2024photoreceptor} showed that sensor layout changes downstream task performance. Neither tests whether non-uniform input reshapes learned representations in navigation agents---our central question.

\textbf{Probing neural representations.}\quad
Belinkov~\cite{belinkov2022probing} surveys probing classifiers for neural networks. Hewitt \& Liang~\cite{hewitt2019control} introduced control tasks to distinguish genuine structure from probe memorization. We adopt both linear probing and the Hewitt--Liang selectivity metric to validate our findings.

% ==========================================================================
% 3. METHODS
% ==========================================================================
\section{Methods}

\subsection{Agent Architecture and Training}

All agents are trained on PointGoal navigation in Habitat~\cite{savva2019habitat} using DD-PPO~\cite{wijmans2020ddppo} on a merged pool of Gibson-0+~\cite{xia2018gibson} (411 scenes) and Matterport3D train (61 scenes), totalling 472 training scenes. Evaluation uses 18 held-out MP3D test scenes (1008 episodes).

All agents share an identical 3-layer LSTM (512-d) backbone and the Wijmans non-visual sensor stack: goal-in-start-frame, GPS, compass, and a close-to-goal scalar, each projected to 32-d and concatenated with a 32-d previous-action embedding. Training targets are set per condition to ensure convergence rather than a fixed frame budget: uniform trains to 400M frames (its ResNet-18 encoder requires more samples), while other conditions target 250M frames.

We observed silent weight corruption in a pilot DD-PPO run caused by NaN gradients flowing through \texttt{clip\_grad\_norm\_} into \texttt{optimizer.step()}; we mitigate this with a gradient-sanitisation patch to \texttt{PPO.before\_step}, described in Appendix~\ref{app:training-stability}.

\subsection{Visual Conditions}

Five conditions vary only the visual input:

\begin{enumerate}
    \item \textbf{Blind:} No visual input. Receives only the non-visual sensor stack, replicating Wijmans et al.~\cite{wijmans2023emergence}.
    \item \textbf{Uniform:} Full-resolution 256$\times$256 RGB encoded by ResNet-18.
    \item \textbf{Foveated (fixed):} Eccentricity-dependent Gaussian blur ($\sigma_{\max}=8$, quadratic falloff) with gaze locked at image centre. Same ResNet-18 encoder.
    \item \textbf{Matched-compute:} Low-resolution uniform RGB that exceeds the foveated effective pixel budget, controlling for total information volume.
    \item \textbf{Foveated (learned):} Same blur model, but gaze centre predicted by a lightweight MLP trained end-to-end with the navigation policy.
\end{enumerate}

The matched-compute condition is critical: by providing more total information than foveation but distributing it uniformly, it isolates the effect of \emph{spatial non-uniformity} from reduced information volume.

\subsection{Probing Methodology}

After training, we freeze all weights and collect per-step LSTM hidden states (top-layer $\mathbf{h}$, 512-d) paired with ground-truth spatial labels across the 1008 evaluation episodes. We train Ridge regression probes ($\alpha=10$) with \textbf{episode-level} train/test splits to prevent information leakage within episodes. Probe quality is validated with the Hewitt--Liang selectivity metric~\cite{hewitt2019control}: we compare probe accuracy on the true task to accuracy on a control task (random permutation of labels), confirming that high $R^2$ reflects genuine spatial structure rather than probe memorization.

Probed variables include: GPS position, compass heading, distance-to-goal, path history (position at lag $k \in \{0,\ldots,5\}$ from current hidden state), visited-region occupancy, full occupancy grid (20$\times$20), and per-unit spatial information (place-cell analysis).

\subsection{Cross-Condition Analysis}

To test H2 (representational divergence), we employ two complementary methods:

\textbf{Linear CKA}~\cite{kornblith2019cka} measures representational similarity between conditions by comparing the geometry of hidden-state activations on matched episodes. Values near 1.0 indicate identical geometry; near 0.0 indicates unrelated representations.

\textbf{Probe transfer:} We train a GPS probe on one condition's hidden states and evaluate it on another's. If spatial information is encoded in a shared linear subspace, cross-condition transfer should succeed. Catastrophic failure indicates fundamentally different encoding formats.

\subsection{Behavioral Tests}

Probes tell us what information is present in memory; they do not tell us whether that information is functionally used. We therefore add a behavioral test inspired by the shortcut-discovery protocol in the SPACE benchmark~\cite{ramakrishnan2025space}, itself a modern analogue of Tolman's classic cognitive-map experiments~\cite{tolman1948cognitive}. For each evaluation scene we run its episodes sequentially in two conditions: (i)~\textbf{persistent} memory, in which the LSTM hidden state carries over between episodes, so the agent accumulates spatial experience within a scene; and (ii)~\textbf{reset} memory, in which the hidden state is zeroed before each episode (the default evaluation protocol). For each condition we measure Success-weighted Path Length (SPL). The persistence benefit $\Delta\text{SPL} = \text{SPL}_\text{persistent} - \text{SPL}_\text{reset}$ quantifies how much accumulated spatial memory improves navigation efficiency. Comparing $\Delta\text{SPL}$ across conditions isolates which agents build the most operationally useful cognitive maps, independent of whether their probes decode well.

% ==========================================================================
% 4. RESULTS
% ==========================================================================
\section{Results}

\subsection{Training Performance}

\begin{table}[t]
\centering
\caption{Task performance per condition on rolling 50-episode training windows (SPL and success reported on the final update; Gibson-0+/MP3D train pool, MP3D test held out). Foveated (resume) and matched-128 runs are in progress at time of writing; we report the partially-trained probed checkpoint.}
\label{tab:training}
\begin{tabular}{lcccl}
\toprule
Condition & Frames & Success & SPL & Status \\
\midrule
Blind            & 342M & 0.928 & 0.588 & Converged \\
Uniform          & 250M & \textbf{0.986} & \textbf{0.852} & Converged \\
Foveated (fixed) & 174M$^\ast$ & --- & --- & Restarted (corrupt)$^\ast$ \\
Matched-48 (deg.) & 250M & 0.979 & 0.670 & Converged, degenerate \\
Matched-128      & 82M  & --- & --- & In progress \\
Fov-learned      & 250M & 0.978 & 0.820 & Converged \\
\bottomrule
\end{tabular}
\end{table}

Table~\ref{tab:training} summarises task performance at the probed checkpoint for each condition. Three conditions (blind, uniform, fov-learned) have fully converged. Matched-48 reached 250M frames but with a degenerate $1\times 1$ feature map (Appendix~\ref{app:training-stability} documents the architectural issue and the matched-128 replacement currently training). The fixed-gaze foveated run initially completed 250M frames but its weights were silently corrupted (Appendix~\ref{app:training-stability}); we are re-training from its last clean intermediate checkpoint (ckpt.36 at 174M frames) and report probing results from that clean snapshot in the remainder of §4.

Among fully converged conditions, uniform achieves the highest task performance (SPL $0.852$, $98.6\%$ success), reflecting the extreme informational advantage of a full-resolution RGB view on short-horizon PointGoal navigation. The fov-learned condition performs comparably ($0.820$ SPL, $97.8\%$), indicating that adding a blur-based foveation and a learned gaze decoder does not substantially degrade task ability. The blind baseline, which solves the task from proprioception alone, reaches $0.588$ SPL at 342M frames, reproducing the qualitative finding of Wijmans et al.~\cite{wijmans2023emergence}. The crucial observation for the representational analysis below is that \emph{all} converged sighted conditions surpass 97\% success, so any differences in probe decoding cannot be attributed to differential task competence.

\subsection{Baseline Probes: GPS, Compass, Distance-to-Goal}

\begin{table}[t]
\centering
\caption{Global linear probe accuracy (Ridge, $\alpha=10$, episode-level split). Foveated-learned values reported on a matched-distribution subset (first 4 steps per episode); see §4.3 caveat.}
\label{tab:probe_baseline}
\begin{tabular}{lcccc}
\toprule
 & GPS $R^2$ & GPS MAE & Compass $R^2$ & Compass MAE \\
\midrule
Blind       & \textbf{0.899} & 0.020m & 0.878 & 0.8\textdegree \\
Foveated    & 0.841          & 0.034m & 0.715 & 2.1\textdegree \\
Matched-48  & 0.832          & 0.033m & 0.607 & 1.3\textdegree \\
Uniform     & 0.614          & 0.035m & 0.423 & 1.6\textdegree \\
Fov-learned & 0.031          & 0.022m & \textbf{0.941} & --- \\
\bottomrule
\end{tabular}
\end{table}

\begin{table}[t]
\centering
\caption{Distance-to-goal probe and Hewitt--Liang selectivity.}
\label{tab:probe_control}
\begin{tabular}{lccc}
\toprule
 & Dist-to-goal $R^2$ & GPS selectivity & Compass selectivity \\
\midrule
Blind    & \textbf{0.984} & \textbf{0.911} & \textbf{0.890} \\
Foveated & 0.953          & 0.880          & 0.751 \\
Matched  & 0.977          & 0.871          & 0.701 \\
Uniform  & 0.959          & 0.763          & 0.543 \\
\bottomrule
\end{tabular}
\end{table}

Table~\ref{tab:probe_baseline} and Figure~\ref{fig:global_probe} show baseline probe results. Blind dominates GPS and compass probes because these signals are directly in its input, so the LSTM trivially encodes them---replicating Wijmans et al.~\cite{wijmans2023emergence}. Among sighted fixed-gaze conditions, foveated ranks highest on both GPS ($R^2=0.841$) and compass ($R^2=0.715$); uniform is weakest on GPS ($R^2=0.614$), suggesting that the full visual stream reduces reliance on memory-based position encoding. Foveated-learned is qualitatively different: it achieves the \emph{highest} compass $R^2$ of all five conditions ($0.941$) yet near-chance GPS ($R^2=0.031$). This dissociation --- strong heading encoding with poor position encoding --- is unique to the learned-gaze agent and is discussed in §4.3 and §4.5.

Table~\ref{tab:probe_control} confirms probe validity: all selectivity values are well above zero (blind $0.91$, foveated $0.88$, matched $0.87$, uniform $0.76$), indicating that probes capture genuine spatial structure rather than memorisation artefacts. Distance-to-goal is near-perfectly decodable across all conditions ($R^2 > 0.95$).

\begin{figure}[t]
\centering
\includegraphics[width=0.95\linewidth]{fig/h1_global_probe.pdf}
\caption{Global GPS probe $R^2$ (solid) and Hewitt--Liang selectivity (hatched) per condition. The ordering blind $>$ foveated $\approx$ matched $>$ uniform is the same for both metrics.}
\label{fig:global_probe}
\end{figure}

\subsection{H1: Compensatory Memory --- Path-History Probe}

\begin{figure}[t]
\centering
\includegraphics[width=0.95\linewidth]{fig/h1_path_history.pdf}
\caption{Path-history probe $R^2$ vs.\ lag $k$. The foveated agent retains decodable spatial information far into the past (lag-5 $R^2=0.556$); uniform decays catastrophically (lag-5 $R^2=-1.53$). Decoding is a Ridge regression ($\alpha=10$) with episode-level train/test splits; negative values indicate the probe fails to generalise.}
\label{fig:path_history}
\end{figure}

\begin{table}[t]
\centering
\caption{Path-history probe: $R^2$ for decoding GPS at lag $k$ from the current hidden state. Higher retention at large $k$ indicates longer working spatial memory. The foveated-learned row will be added once its re-analysis completes.}
\label{tab:path_history}
\begin{tabular}{lcccccc}
\toprule
 & lag-0 & lag-1 & lag-2 & lag-3 & lag-4 & lag-5 \\
\midrule
Blind    & 0.899 & 0.905 & 0.867 & 0.748 & 0.536 & 0.043 \\
Foveated & 0.841 & 0.853 & 0.824 & 0.737 & \textbf{0.629} & \textbf{0.556} \\
Matched  & 0.832 & 0.855 & 0.807 & 0.581 & 0.445 & 0.274 \\
Uniform  & 0.614 & 0.595 & 0.534 & 0.386 & $-$0.105 & $-$1.532 \\
\bottomrule
\end{tabular}
\end{table}

The path-history probe (Figure~\ref{fig:path_history}, Table~\ref{tab:path_history}) is the paper's strongest result. This probe decodes the agent's GPS position $k$ steps in the past from the current hidden state, testing how far back working spatial memory extends.

At lag-5, the foveated agent achieves $R^2=0.556$, compared with matched-compute ($R^2=0.274$), blind ($R^2=0.043$), and uniform ($R^2=-1.532$; probe fails to generalise). Retention relative to lag-0 is 66\% for foveated, 33\% for matched, 5\% for blind, and negative for uniform. Foveated is the only condition that retains more than half of its immediate-past encoding quality at 5 steps back.

Two comparisons isolate the mechanism. (i)~foveated~$>$~uniform at every lag~$k$: although both observe the environment, the uniform agent's full-resolution visual stream removes the \emph{necessity} of remembering where it has been. (ii)~foveated~$>$~matched at every lag~$k \geq 3$: matched-compute provides more \emph{total} visual pixels than foveated but distributes them uniformly, so total information volume does not drive the effect. It is the \emph{spatial non-uniformity} of the foveated sensor that forces compensatory temporal memory.

A related timing view (Figure in Appendix, same data) shows that the spatial encoding rises steeply over the first three timesteps of each episode and stabilises thereafter; the compensatory effect is thus present throughout an episode rather than being a late-emergent behaviour.

\paragraph{Foveated-learned does not fit the retention picture.} On the same probe, the foveated-learned agent shows essentially flat $R^2\approx-2.5$ across all lags $k=0,\dots,5$ (Figure~\ref{fig:path_history}, orange). This is not a retention failure of degree but a qualitative difference in what is encoded: §4.2 shows fov-learned has the \emph{highest} compass decoding of all conditions ($R^2=0.941$) but near-chance GPS decoding ($R^2=0.031$), when evaluated on a matched-distribution subset. In other words, the foveated-learned hidden state carries heading information strongly but absolute-position information weakly --- inconsistent with the compensatory-memory story seen for fixed-gaze foveation. We return to this in §4.5.

\subsection{H2: Representational Divergence --- CKA and Probe Transfer}

\begin{figure}[t]
\centering
\includegraphics[width=0.8\linewidth]{fig/h2_cka_heatmap.pdf}
\caption{Unaligned linear CKA~\cite{kornblith2019cka} between top-layer LSTM hidden states of every pair of conditions (truncated to common $n=1{,}907$). Every off-diagonal value is below $0.004$; the most-similar pair is foveated-fixed vs.\ foveated-learned ($0.004$). Colour scale clipped at 0.05.}
\label{fig:cka}
\end{figure}

\begin{table}[t]
\centering
\caption{Unaligned linear CKA on top-layer hidden states ($n=1{,}907$ per condition). All pairs of conditions develop nearly-orthogonal representational geometries.}
\label{tab:cka}
\begin{tabular}{lccccc}
\toprule
 & Blind & Uniform & Fov-fix & Fov-lrn & Matched \\
\midrule
Blind    & ---   & 0.0010 & 0.0016 & 0.0011 & 0.0026 \\
Uniform  &       & ---    & 0.0011 & 0.0019 & 0.0012 \\
Fov-fix  &       &        & ---    & \textbf{0.0035} & 0.0008 \\
Fov-lrn  &       &        &        & ---    & 0.0014 \\
\bottomrule
\end{tabular}
\end{table}

\begin{table}[t]
\centering
\caption{Probe-transfer $R^2$ on GPS: probe trained on row condition, tested on column condition. Self-accuracy (diagonal) is $0.90$--$0.97$; every cross-condition transfer is catastrophically negative. Foveated-learned shows exceptionally large cross-transfer deficits (see text).}
\label{tab:transfer}
\begin{tabular}{lccccc}
\toprule
Train $\backslash$ Test & Blind & Uniform & Fov-fix & Fov-lrn & Matched \\
\midrule
Blind    & \textbf{0.97} & $-$9.7   & $-$2.7   & $-$613   & $-$15.7 \\
Uniform  & $-$7.3        & \textbf{0.97} & $-$13.5 & $-$410  & $-$11.5 \\
Fov-fix  & $-$11.3       & $-$5.1   & \textbf{0.96} & $-$232  & $-$10.7 \\
Fov-lrn  & $-$482        & $-$430   & $-$456   & \textbf{0.90} & $-$855 \\
Matched  & $-$5.7        & $-$11.7  & $-$15.5  & $-$390   & \textbf{0.96} \\
\bottomrule
\end{tabular}
\end{table}

Table~\ref{tab:cka} and Figure~\ref{fig:cka} show that every cross-condition CKA value is below $0.004$. No two conditions encode space with similar geometry. Blind is the most dissimilar from sighted conditions, consistent with its reliance on proprioceptive rather than visual features. The closest pair is the two foveated variants (fixed-gaze vs.\ learned-gaze), yet even this pair is only at $0.0035$.

Table~\ref{tab:transfer} confirms the divergence behaviourally: cross-condition probe transfer is catastrophically negative for every non-self pair (self $R^2 = 0.90$--$0.97$; cross $R^2 \leq -2.7$). Spatial information is not encoded in a shared linear subspace across conditions.

\paragraph{Caveat on foveated-learned.} The foveated-learned row and column in Table~\ref{tab:transfer} show cross-transfer values that are one-to-two orders of magnitude worse than other pairs ($-232$ to $-855$ rather than $-2$ to $-15$). This is dominated by a probing-protocol mismatch rather than by representational content: the foveated-learned agent's evaluation episodes averaged 171 steps per episode (median length 77), vs.\ $\approx 4$ steps per episode for the other conditions, so its probe dataset covers a spatial range with variance $19$\,m$^2$ vs.\ $\sim 10^{-4}$\,m$^2$ for the others. On a truncated-to-matched subset (first 4 steps of each episode, $n\approx 2000$), cross-condition transfer values fall into the same $[-2, -15]$ band as the other pairs (reported in Appendix~\ref{app:h2_matched_subset}, job 2841522). The qualitative conclusion --- zero cross-condition sharing of spatial code --- is unchanged.

Combined, near-zero CKA across all pairs and catastrophic probe-transfer failure provide strong evidence for H2: visual input type fundamentally shapes the representational geometry used to encode space, not just the amount of spatial information stored.

\subsection{H3: Epistemic Gaze --- Preliminary Observations}

H3 predicts that when gaze is learned end-to-end through the navigation loss, it will target regions where memory is most uncertain, amplifying the H1--H2 effects relative to fixed-centre foveation. Testing this requires a learned-gaze agent with \emph{non-trivial} gaze behaviour; our implementation did not produce one.

\paragraph{Observation 1: gaze collapse.} After 250M frames of training, the foveated-learned agent's gaze decoder converges to a near-constant output. Across 85\,K evaluation steps, the per-dimension standard deviation is $5\times 10^{-4}$ and $2\times 10^{-4}$ on the two gaze coordinates --- an effective spatial range of $<1\%$ of the image. Functionally, the learned-gaze agent behaves as a fixed-gaze agent at an offset position $(0.49, 0.62)$ rather than $(0.5, 0.5)$.

\paragraph{Observation 2: different representational content.} Despite the gaze collapse, the agent reaches $0.820$ SPL and $97.8\%$ success (Table~\ref{tab:training}) --- comparable to the fixed-gaze foveated agent. Its representation, however, is qualitatively different: it has the highest compass decoding of all five conditions ($R^2=0.941$, Table~\ref{tab:probe_baseline}) but near-chance GPS decoding ($R^2=0.031$). In representational-similarity space, it is the closest of all pairs to fixed-gaze foveated (CKA $0.0035$, Figure~\ref{fig:cka}), but probe weights learned on fixed-gaze fail to transfer ($R^2=-232$, Table~\ref{tab:transfer}).

\paragraph{Interpretation and caveats.} We interpret these as \emph{preliminary observations}, not a refutation of H3. With a single seed, a single gaze-decoder architecture (sigmoid MLP from a detached hidden state), and a single rollout approximation (``slow-gaze'', one gaze per 256-step segment), we cannot determine whether gaze collapse is a fundamental consequence of implicit task pressure, an artefact of sigmoid saturation blocking gradient flow, or a single-seed local minimum. A gaze-diversity auxiliary loss (hinge penalty on sub-target variance, Appendix B) is in pilot evaluation; if it yields non-trivial gaze without destroying task performance, we will re-run H1--H3 analyses with that agent. What we can state firmly is that \emph{in the setup we tested}, the implicit navigation signal is not sufficient to produce actively moving gaze, and the resulting policy encodes heading rather than position.

\subsection{Additional Analyses}

\textbf{Visited-region probe.}\quad
All conditions track visited regions near-perfectly ($>$99.8\% binary accuracy; $R^2 \in [0.801, 0.834]$), suggesting this is a basic capability shared by all agents regardless of visual input.

\textbf{Occupancy map reconstruction.}\quad
Occupancy decoding from a 20$\times$20 grid is weak across all conditions (all $R^2 < 0$; binary accuracy 55--60\%). None of the agents build detailed layout maps in their hidden states at this resolution. This may require a finer grid, more evaluation episodes, or a non-linear decoder.

\textbf{Place-cell-like units.}\quad
Nearly all LSTM units carry spatial information ($>$1 bit) across all conditions: foveated 475/512, blind 473/512, uniform 473/512, matched 461/512. Spatial encoding is broadly distributed rather than concentrated in a small population.

\textbf{Multi-layer comparison.}\quad
For the blind agent, spatial information is distributed across LSTM layers but concentrated in layers 0--1, with layer-0 hidden state achieving the highest GPS $R^2$ (0.958).

\textcolor{red}{TODO: Add shortcut discovery behavioral test results once analysis is complete. Add PCA/t-SNE visualization of hidden states colored by position and condition.}

% ==========================================================================
% 5. DISCUSSION
% ==========================================================================
\section{Discussion}

\subsection{Sensor Structure as a Representational Axis}

Our results establish that sensor structure---specifically, the spatial distribution of visual quality---causally shapes what recurrent memory learns to represent during navigation. The foveated agent does not simply store degraded versions of the uniform agent's representations; it builds qualitatively different spatial codes with stronger temporal memory (H1) embedded in a fundamentally different geometry (H2).

The path-history result is particularly striking: the foveated agent retains position information over longer horizons than even the blind agent, which receives GPS directly. This suggests that spatially non-uniform input creates a pressure to maintain richer egocentric path traces---a form of compensatory memory that arises from pure RL without explicit probabilistic objectives or architectural priors for uncertainty, consistent with the belief-state emergence reported by Hennig et al.~\cite{hennig2023belief}.

\subsection{The Role of Spatial Non-Uniformity}

The matched-compute control is essential for interpretation. Matched-compute provides more total visual information than foveation but distributes it uniformly, yet its path-history retention ($R^2=0.445$ at lag-5) falls well below foveated ($R^2=0.629$). This isolates the causal factor: it is the spatial non-uniformity of visual quality, not merely reduced information volume, that drives compensatory memory. This is an information-bottleneck effect driven by sensor structure rather than architectural choice.

\subsection{Why Uniform Agents Have Weak Memory}

Uniform vision yields the weakest spatial encoding across nearly all probes. We interpret this as a form of representational outsourcing: when current visual input is reliable, the agent can re-derive spatial information from each frame rather than maintaining it in memory. This is efficient for the task but produces shallow temporal traces. Foveated and blind agents cannot rely on current visual input and must therefore invest in memory.

\subsection{Limitations}

Several limitations constrain the generality of our findings. First, all results use intermediate checkpoints; final-checkpoint probing may shift quantitative conclusions. Second, the foveated-learned condition (H3) has not yet been trained, leaving the epistemic gaze hypothesis untested. Third, our foveation model (Gaussian blur with quadratic falloff) is a simplification of biological foveation; log-polar or multi-scale pyramid models may produce different effects. Fourth, we probe only linear decoders---non-linear probes might reveal spatial structure that Ridge regression misses, particularly for occupancy maps. Fifth, the LSTM backbone limits generality to transformer-based architectures.

\subsection{Broader Implications}

If sensor design predictably shapes the structure of learned spatial codes, it becomes a tunable axis---alongside architecture and objective---for steering embodied agents toward representations that are more interpretable, more robust, or more aligned with human spatial cognition. The finding that degraded input can produce \emph{stronger} memory traces (rather than merely weaker ones) challenges the assumption that richer input always yields richer representations and suggests that controlled information bottlenecks may be a design tool for representational engineering.

% ==========================================================================
% 6. CONCLUSION
% ==========================================================================
\section{Conclusion}

We investigated how foveated vision shapes cognitive maps in navigation agents by training five conditions that vary only in visual input and probing their LSTM hidden states with linear decoders. Our two main findings are: (1)~foveated vision induces compensatory temporal memory---the foveated agent retains spatial information over longer horizons than all other conditions, including blind agents with direct GPS access; and (2)~each visual condition produces a fundamentally different representational geometry (CKA~$<$~0.003, catastrophically negative probe transfer), despite comparable task performance. These results establish sensor structure as a causal factor in the structure of learned spatial representations, complementing prior work on architecture and objective design.

\textcolor{red}{TODO: Update conclusion with H3 results and behavioral test findings once available.}

\bibliographystyle{IEEEtran}
\bibliography{literature}

\appendices

\section{Training Stability: Silent Weight Corruption in DD-PPO}
\label{app:training-stability}

During an initial foveated training run we observed a failure mode that did not manifest as a crash but as slow, silent corruption of the learned weights. After roughly 170M frames of stable progress (reward~$\approx$~10, SPL~$\approx$~0.83), one mini-batch on one DDP worker produced a NaN gradient. Because \texttt{torch.nn.utils.clip\_grad\_norm\_} does not sanitise non-finite gradients---the global norm becomes NaN, the clip coefficient becomes NaN, and \texttt{grad *= NaN} remains NaN---the subsequent \texttt{optimizer.step()} wrote NaN into every parameter. DDP all-reduce then propagated the corruption across workers. The training loop did not raise: it continued for 2.5 days with reward~$=$~NaN and SPL~$=$~0, until \texttt{torch.multinomial} finally rejected the NaN probability tensor. All checkpoints saved during this window were unrecoverable (90~of 97~parameter tensors entirely NaN).

To isolate the origin of the NaN gradient we re-ran training from the last clean checkpoint with \texttt{torch.autograd.set\_detect\_anomaly(True)} and per-module forward-output hooks. The anomaly traceback pinned the failure to \texttt{MseLossBackward0}, i.e.\ the backward pass of the value-head loss~$0.5 \cdot (V(s) - R)^2$; rare return-vs-value discrepancies overflowed float32 in the squaring step, producing non-finite gradients downstream.

Our mitigation is a minimally invasive patch to \texttt{habitat\_baselines.rl.ppo.ppo.PPO.before\_step} that zeroes any non-finite gradient element before clipping (\texttt{torch.nan\_to\_num\_} with \texttt{nan=posinf=neginf=0}). A mini-batch with any NaN gradient thus becomes a no-op update rather than a catastrophic one. The patch is installed at import time when \texttt{src.habitat} is loaded and is a no-op on clean training paths (\texttt{nan\_to\_num} of a finite tensor is identity). Across the five conditions reported in this paper, no NaN events were observed during full training after the patch was deployed, and all final checkpoints contain only finite weights. We reported the underlying habitat-baselines issue upstream (\texttt{facebookresearch/habitat-lab} \#2226) with a minimal reproduction and a proposed native fix.

\end{document}
